% ============ 第2章 §2.1 Lebesgue 测度 (原书 pp.36-43) ============

第1章通过数列概观了调和分析, 但一般地, 调和分析中大量使用积分. 又, 由于必须基于完备性进行讨论, 几乎都使用 Banach 空间. 本章叙述 Lebesgue 积分的基础事项以及函数解析的基础事项. 关于 Lebesgue 积分的基础事项, 不作证明地总结. 详细证明请参考适当的文献. 又, 这一章几乎没有练习题.

\section{Lebesgue 测度}\label{sec:2.1}
\index{0symbolcap@$\bigcap_{\lambda \in \Lambda} U_{\lambda}$}

\index{0symbolcup@$\bigcup_{\lambda \in \Lambda} U_{\lambda}$}

\index{0symbolBxr@$B(x, r)$}

\index{0symbolsupp@$\operatorname{supp}(f)$}

\index{Z@支集}

\index{J@简单函数}

\index{C@长方体}

\index{N@内正则性}

\index{L@Lebesgue 外测度}

\index{L@Lebesgue 可测函数}

\index{L@Lebesgue 可测集合}
























这里, 除证明外, 简洁地总结 Lebesgue 积分的基础事项.

\subsection*{集合的并, 交集}

调和分析亦可说是研究函数性质的学问, 在调和分析中, 集合的操作是必需事项. 本书使用以下记号.

\begin{definition}\label{def:2.1}
(1) 集合的子族, 亦即子集族 $\{U_{\lambda}\}_{\lambda\in\Lambda}\subset\mathbb{R}^{n}$, 是指

(A) $\Lambda$ 是某个集合, 且

(B) 对 $\lambda \in \Lambda$, 给定 $\mathbb{R}^{n}$ 的子集 $U_{\lambda}$.

此时, 称 $\{U_{\lambda}\}_{\lambda\in\Lambda}$ 为 $\mathbb{R}^{n}$ 的子集族.

(2) 对 $\mathbb{R}^{n}$ 的子集族 $\{U_{\lambda}\}_{\lambda\in\Lambda}\subset\mathbb{R}^{n}$, 其并集定义为
\[
\bigcup_{\lambda \in \Lambda}U_{\lambda} = \{x \in \mathbb{R}^{n} : \text{至少存在一个 } \lambda \in \Lambda \text{ 使得 } x \in U_{\lambda}\}.
\]

(3) 对 $\mathbb{R}^{n}$ 的子集族 $\{U_{\lambda}\}_{\lambda\in\Lambda}\subset\mathbb{R}^{n}$, 其交集定义为
\[
\bigcap_{\lambda \in \Lambda}U_{\lambda} = \{x \in \mathbb{R}^{n} : \text{对所有 } \lambda \in \Lambda \text{ 有 } x \in U_{\lambda}\}.
\]
\end{definition}

\subsection*{Lebesgue 可测集合与 Lebesgue 可测函数}

为了研究函数的性质, 时常会用到完备性的讨论, 但如此一来在 Riemann 积分的框架下就不方便了. 这是因为在调和分析中, 几乎所有的情形都是通过由积分定义的范数进行分析的. 因此, 来复习 Lebesgue 可测集合与 Lebesgue 可测函数.

\begin{definition}[直方体]\label{def:2.2}
用满足 $-\infty < a_{j} \leq b_{j} < \infty$ $(j = 1, 2, \ldots, n)$ 的 $2n$ 个实数表为
\[
\prod_{j=1}^{n}[a_{j}, b_{j}) = [a_{1}, b_{1}) \times [a_{2}, b_{2}) \times \cdots \times [a_{n}, b_{n})
\]
的集合称为直方体. 直方体的体积定义为
\[
\left|\prod_{j=1}^{n}[a_{j}, b_{j})\right| = \prod_{j=1}^{n}(b_{j} - a_{j}).
\]
将全体直方体构成的集合记为 $\mathcal{R}$.
\end{definition}

\begin{definition}[Lebesgue 外测度]\label{def:2.3}
设 $A \subset \mathbb{R}^{n}$. $A$ 的 Lebesgue 外测度 $m^{*}(A)$ 是由
\[
m^{*}(A) = \Gamma(A) \equiv \inf\left\{\sum_{j=1}^{\infty}|R_{j}| : \text{各 } R_{j} \text{ 是直方体且 } A \subset \bigcup_{j=1}^{\infty}R_{j}\right\}
\]
给出的属于 $[0, \infty]$ 的量.
\end{definition}

\begin{definition}[Lebesgue 可测集合]\label{def:2.4}
当对 $\mathbb{R}^{n}$ 的所有子集 $F$ 都有 $\Gamma(F) = \Gamma(F \cap E) + \Gamma(F \cap E^{c})$ 成立时, 称 $E$ 为 Lebesgue 可测集合, 或 Lebesgue 集合. 此时规定 $|E| = \Gamma(E)$.
\end{definition}

在 Lebesgue 测度的性质中, 称为内正则性的性质很方便, 复习一下这个性质.

\begin{theorem}[Lebesgue 测度的内正则性]\label{thm:2.1}
Lebesgue 测度的内正则性是指, 对任意可测集合 $A$,
\begin{equation}\tag{2.1}
|A| = \sup\{|K| : K \text{ 是包含于 } A \text{ 的紧集}\}
\label{eq:2.1}
\end{equation}
成立.
\end{theorem}

\begin{proof}
这里不作证明. 请参考适当的文献.
\end{proof}

\begin{definition}\label{def:2.5}
设 $E$ 为 Lebesgue 可测集合. $f: E \to \overline{\mathbb{R}} = \mathbb{R} \cup \{\pm\infty\}$ 是 Lebesgue 可测函数, 是指对任意 $\lambda \in \mathbb{R}$,
\[
f^{-1}(\lambda, \infty] = \{x \in E : f(x) > \lambda\}
\]
是 Lebesgue 可测集合.
\end{definition}

\subsection*{非负值 Lebesgue 可测函数积分的定义}

比区间上的积分更一般地, 考虑集合 $E$ 上的积分. 简单函数是指只取有限个值的可测函数, 且在某个直方体之外为 $0$ 的函数.

\begin{definition}[简单函数的积分]\label{def:2.6}
设 $E$ 为 Lebesgue 可测集合, $f: E \to [0, \infty)$ 为非负值简单函数. 用 $a_{1}, a_{2}, \ldots, a_{N} \in \mathbb{R}$ 与 (不允许有公共部分的) Lebesgue 可测集合 $E_{1}, E_{2}, \ldots, E_{N}$ 表为 $f = \sum_{j=1}^{N}a_{j}\chi_{E_{j}}$ 时,
\[
\int_{E}f(x)\dd x = \sum_{j=1}^{N}a_{j}|E \cap E_{j}| = \sum_{j=1}^{N}a_{j}|E_{j}|
\]
成立.
\end{definition}

非负值函数是指取值于 $[0, \infty]$ 的函数.

\begin{definition}[非负值函数的积分]\label{def:2.7}
设 $E$ 为 Lebesgue 可测集合. 当 $f: E \to [0, \infty]$ 为非负值 Lebesgue 可测函数时,
\[
\int_{E}f(x)\dd x \equiv \sup\left\{\int_{E}g(x)\dd x : g \text{ 是简单函数且对所有 } x \in E \text{ 有 } 0 \leq g(x) \leq f(x)\right\} \quad (\geq 0)
\]
成立. $\int_{E}f(x)\dd x = \infty$ 的可能性也存在.
\end{definition}

\subsection*{单调收敛定理, Fatou 引理}

可测函数这个说法是笼统的, 用容易理解的函数来逼近是调和分析中也常出现的手法之一. 下面的定理经常使用.

\begin{theorem}[单调收敛定理]\label{thm:2.2}
设 $\{f_{j}\}_{j=1}^{\infty}$ 是在 Lebesgue 可测集合 $E$ 上定义的非负值 Lebesgue 可测函数的增列. 用 $\lim_{j\to\infty}f_{j}(x) = f(x)$ $(x \in E)$ 定义函数 $f$ 时,
\begin{equation}\tag{2.2}
\lim_{j \to \infty}\int_{E}f_{j}(x)\dd x = \int_{E}f(x)\dd x
\label{eq:2.2}
\end{equation}
成立.
\end{theorem}

\begin{proof}
这里不作证明. 请参考适当的文献.
\end{proof}

\begin{theorem}[Fatou 引理]\label{thm:2.3}
对在 Lebesgue 可测集合 $E$ 上定义的非负值 Lebesgue 可测函数列 $\{f_{j}\}_{j=1}^{\infty}$,
\begin{equation}\tag{2.3}
\int_{E}\liminf_{j \to \infty}f_{j}(x)\dd x \leq \liminf_{j \to \infty}\int_{E}f_{j}(x)\dd x
\label{eq:2.3}
\end{equation}
成立.
\end{theorem}

\begin{proof}
这里不作证明. 请参考适当的文献.
\end{proof}

给出 Fatou 引理的一个变形. 对函数 $f: (a, b) \to \mathbb{R}$ 与 $t_{0} \in [a, b]$,
\[
\liminf_{u \to t_{0}, u \in (a, b)}f(u) = \sup_{\varepsilon > 0}\left(\inf_{u \in (a, b) \cap (t_{0} - \varepsilon, t_{0} + \varepsilon)}f(u)\right)
\]
成立.

\begin{theorem}\label{thm:2.4}
设 $E$ 为 Lebesgue 可测集合. $F: E \times (a, b) \to [0, \infty]$ 是对每个 $t \in (a, b)$, $F(\cdot, t)$ 与 $\liminf_{u \to t_{0}, u \in (a, b)}F(\cdot, u)$ 成为 Lebesgue 可测的函数. 此时, 对任意 $a \leq t_{0} \leq b$,
\[
\liminf_{t \to t_{0}, t \in (a, b)}\int_{E}F(x, t)\dd x \geq \int_{E}\liminf_{t \to t_{0}}F(x, t)\dd x
\]
成立. (这里, $t_{0} = b$, $b = \infty$ 也可以.)
\end{theorem}

\begin{proof}
这里不作证明. 请参考适当的文献.
\end{proof}

\subsection*{Lebesgue 可积函数}

作为调和分析的课题之一, 可举出函数各种可积性的测量. 这里暂且引入绝对可积的定义. 设 $E$ 为 Lebesgue 可测集合. Lebesgue 可测函数 $f: E \to [-\infty, \infty]$ 可积与 $|f|$ 可积等价.

$L^{1}(\mathbb{R}^{n})$ 表示满足
\[
\|f\|_{L^{1}} = \int_{\mathbb{R}^{n}}|f(x)|\dd x < \infty
\]
的可测函数 $f$ 全体构成的集合. 又, $L^{\infty}(\mathbb{R}^{n})$ 表示满足
\[
\|f\|_{L^{\infty}} = \inf\{\lambda > 0 : \text{在测度 } 0 \text{ 的集合 } N \text{ 之外, } |f(x)| < \lambda\} < \infty
\]
的全体构成的集合. 这里, $\|f\|_{L^{1}}$, $\|f\|_{L^{\infty}}$ 分别称为 $L^{1}$-范数, $L^{\infty}$-范数. 不过, 只在测度 $0$ 的集合上不同的函数视为同一. 这些范数在本书中经常使用. 记号 $\lesssim$ 是指 $\leq c$, 这里的 $c$ 在同一的定理、证明等中, 多个值也使用同一个字母. 例如 $3c \leq c$ 等就是其典型用例.

\begin{example}\label{exa:2.1}
$F(x) = (1 + |x|)^{-n-1}$, $x \in \mathbb{R}^{n}$ 是可积函数. 即 $F \in L^{1}(\mathbb{R}^{n})$. 实际上, 由单调收敛定理 (定理~\ref{thm:2.2}),
\[
\int_{\mathbb{R}^{n}}F(x)\dd x = \int_{|x| < 1}F(x)\dd x + \sum_{j=1}^{\infty}\int_{2^{j-1} < |x| < 2^{j}}F(x)\dd x
\]
成立, 但
\begin{align*}
\int_{2^{j-1} < |x| < 2^{j}}F(x)\dd x &\leq (1 + 2^{j-1})^{-n-1}\int_{2^{j-1} < |x| < 2^{j}}1\dd x \\
&\lesssim 2^{jn}(1 + 2^{j-1})^{-n-1} \\
&\leq 2^{jn-(j-1)(n+1)}
\end{align*}
成立, 所以把它关于 $j$ 相加即可.
\end{example}

\subsection*{Lebesgue 收敛定理及其变形}

\begin{theorem}[Lebesgue 收敛定理]\label{thm:2.5}
设 $E$ 为 Lebesgue 可测集合. 设 $f_{1}, f_{2}, \ldots: E \to [-\infty, \infty]$ 为逐点收敛到 $f$ 的 Lebesgue 可测函数列. 若存在可积函数 $g: E \to [-\infty, \infty]$ 满足对所有 $j \in \mathbb{N}$ 与 $x \in \mathbb{R}$ 有 $|f_{j}(x)| \leq g(x)$, 则
\begin{equation}\tag{2.4}
\lim_{j \to \infty}\int_{E}f_{j}(x)\dd x = \int_{E}f(x)\dd x
\label{eq:2.4}
\end{equation}
成立.
\end{theorem}

\begin{proof}
这里不作证明. 请参考适当的文献.
\end{proof}

\begin{theorem}\label{thm:2.6}
设 $E$ 为 Lebesgue 可测集合, $F: (a, b) \times E \to \mathbb{R}$ 是满足以下条件的函数.

(1) 对所有 $t \in (a, b)$ 与 $x \in \mathbb{R}$, $\partial_{t}F(t, x)$ 存在.

(2) $F(t, \cdot)$ 可积.

(3) 存在可积函数 $g: E \to \mathbb{R}$, 当 $x \in E$, $t \in (a, b)$ 时, $|\partial_{t}F(t, x)| \leq g(x)$ 成立.

此时,
\begin{equation}\tag{2.5}
\frac{\dd}{\dd t}\int_{E}F(t, x)\dd x = \int_{E}\partial_{t}F(t, x)\dd x
\label{eq:2.5}
\end{equation}
成立.
\end{theorem}

\begin{proof}
这里不作证明. 请参考适当的文献.
\end{proof}

\begin{theorem}\label{thm:2.7}
设 $E$ 为 Lebesgue 可测集合. 设 $F: E \times (a, b) \to \mathbb{R}$ 满足以下条件.

(A) 对所有 $t \in (a, b)$, $\int_{E}|F(x, t)|\dd x < \infty$ 成立.

(B) 对所有 $x \in E$, 函数 $t \in (a, b) \mapsto F(x, t) \in \mathbb{R}$ 连续.

(C) 存在某个可积函数 $G$, 对所有 $t \in (a, b)$ 与 $x \in E$, $|F(x, t)| \leq G(x)$ 成立.

此时,
\[
\lim_{t \to t_{0}}\int_{E}F(x, t)\dd x = \int_{E}F(x, t_{0})\dd x
\]
成立.
\end{theorem}

\begin{proof}
这里不作证明. 请参考适当的文献.
\end{proof}

\begin{theorem}\label{thm:2.8}
设 $E$ 为 Lebesgue 可测集合. 每固定 $t$, $F(\cdot, t)$ 可测的函数 $F: E \times (a, b) \to [0, \infty]$ 满足: 若 $a < t_{1} < t_{2} < b$, 则 $F(x, t_{1}) \leq F(x, t_{2})$. 此时,
\[
\lim_{t \uparrow b}\int_{E}F(x, t)\dd x = \int_{E}F(x, b)\dd x
\]
成立. (这里, $b = \infty$ 也可以.)
\end{theorem}

\begin{proof}
这里不作证明. 请参考适当的文献.
\end{proof}

\subsection*{Fubini 定理}

这里讨论等式
\begin{align*}
(\star): \quad \int_{\mathbb{R}^{m} \times \mathbb{R}^{n}}f(x, y)\dd x\dd y &= \int_{\mathbb{R}^{m}}\left(\int_{\mathbb{R}^{n}}f(x, y)\dd y\right)\dd x \\
&= \int_{\mathbb{R}^{n}}\left(\int_{\mathbb{R}^{m}}f(x, y)\dd x\right)\dd y
\end{align*}
的成立与否. 此后, 括弧写作``等式''时指这个等式. 这里, 关于等式, 也包含了
\[
x \mapsto \int_{\mathbb{R}^{n}}f(x, y)\dd y, \quad y \mapsto \int_{\mathbb{R}^{n}}f(x, y)\dd x
\]
的可测性.

用 $\mathcal{B}$ 表示 $\mathbb{R}^{n}$ 中包含开集的最小 $\sigma$-代数, 称为 Borel $\sigma$-代数. 也就是说, 把满足以下条件的集合族 $\mathcal{Z}$ 全体作为 $\mathcal{S}$:

\begin{itemize}
\item 直方体, $\mathbb{R}^{n}$ 中的空集合 $\emptyset$, $\mathbb{R}^{n}$ 本身属于 $\mathcal{Z}$. 即, 设 $R$ 为任意直方体时, $R, \emptyset, \mathbb{R}^{n} \in \mathcal{Z}$ 成立.

\item 若 $A \in \mathcal{Z}$, 则 $A^{\mathrm{c}} \in \mathcal{Z}$ 成立.

\item 若 $A_{1}, A_{2}, \ldots, A_{j}, \ldots \in \mathcal{Z}$, 则 $\bigcup_{j=1}^{\infty}A_{j} \in \mathcal{Z}$ 成立.
\end{itemize}

由此,
\[
\mathcal{B} = \{B \in 2^{\mathbb{R}^{n}} : \text{对所有 } \mathcal{Z} \in \mathcal{S} \text{ 有 } B \in \mathcal{Z} \text{ 成立}\}
\]
成立. 又, $\mathcal{B}$ 的集合称为 Borel 集合. 由于开集合可表为直方体的可数无限个并, 开集合和闭集合也都是 Borel 集合.

与 Lebesgue 可测函数一样, Borel 可测函数等也可用类似的方法定义. 即, $f: \mathbb{R}^{n} \to \overline{\mathbb{R}}$ 是 Borel 可测函数, 是指 $f^{-1}(\lambda, \infty]$ 是 Borel 可测集合. 例如, 连续函数是 Borel 可测函数. 设 $f$ 为 Borel 可测函数, 对 $f: \mathbb{R}^{m} \times \mathbb{R}^{n} \to [0, \infty]$ 或可积函数 $f: \mathbb{R}^{m} \times \mathbb{R}^{n} \to \mathbb{C}$, $(\star)$ 成立.

\subsection*{Lebesgue 积分与 Riemann 积分的比较}

已述在调和分析中自然地以 Lebesgue 积分的框架考虑积分, 但利用 Riemann 积分的知识来计算积分是重要的. 下面的定理成为这两种积分的桥梁.

\begin{theorem}\label{thm:2.9}
设 $f: [0, \infty) \to \mathbb{R}$ 为连续函数, 假设在广义 Riemann 积分的意义下, 或 Lebesgue 积分的意义下
\[
\int_{0}^{\infty}|f(x)|\dd x < \infty
\]
成立. 此时, 可以说以下事实.

(1) 在一方意义下成立的假设, 实际上在双方意义下都成立. 无论按哪种意义考虑, $\int_{0}^{\infty}|f(x)|\dd x$ 的值都相等.

(2) 在广义 Riemann 积分的意义下, 或 Lebesgue 积分的意义下
\[
\int_{0}^{\infty}f(x)\dd x
\]
存在, 且无论按哪种意义考虑值都相等.
\end{theorem}

\begin{proof}
这里不作证明. 请参考适当的文献.
\end{proof}

\subsection*{函数的支集}

为了不仅调查积分的大小, 还调查函数具有多大的扩展范围, 测量函数的大小是调和分析中的重要主题. 这里, 作为其基准, 说明函数的支集.

\begin{definition}[连续函数的支集]\label{def:2.8}
连续函数 $f: \mathbb{R}^{n} \to \mathbb{C}$ 的支集记作 $\operatorname{supp}(f)$, 定义为 $\{x \in \mathbb{R}^{n} : f(x) \neq 0\}$ 的闭包.
\end{definition}

在几何学教科书中, 把定义~\ref{def:2.8} 作为支集的定义给出, 但 Lebesgue 可测函数的支集也能定义. 此后, 对 $x \in \mathbb{R}^{n}$ 与 $r > 0$, 本书用 $B(x, r)$ 表示以 $x$ 为中心、半径为 $r$ 的开球. $x = 0$ 的情形省略``$x$, '', 记为 $B(r)$.

\begin{definition}[可测函数的支集]\label{def:2.9}
可测函数 $f: \mathbb{R}^{n} \to \mathbb{C}$ 的支集是指, 对任意 $r > 0$, 满足 $|\{y \in B(x, r) : f(y) \neq 0\}| > 0$ 的 $x \in \mathbb{R}^{n}$ 全体.
\end{definition}

对连续函数, 上述两种支集的定义一致. 虽然给出了两个定义, 但基于这一点, 两者都记作 $\operatorname{supp}(f)$. 用 $L_{\mathrm{c}}^{\infty}(\mathbb{R}^{n})$ 表示支集紧的 $L^{\infty}(\mathbb{R}^{n})$ 函数全体. 具有紧支集的函数一般多有可积性, 可以认为能以在有限测度空间上分析的感觉来处理函数. 恒等函数 $1$ 不可积, 而对球 $B(x, r)$, $\chi_{B(x, r)}$ 可积.
